\uuid{dRbq}
\chapitre{Probabilité discrète}
\niveau{L2}
\module{Probabilité}
\sousChapitre{Lois de distributions}
\cpge{mp-mpi.variables-aleatoires-discretes.g, pc.probabilites-lois-usuelles.d, psi.probabilites-lois-usuelles.d, pt.probabilites-lois-usuelles.e, bcpst2.variables-aleatoires-discretes.d}
\titre{Lois usuelles}
\theme{variables aléatoires discrètes, loi géométrique}
\auteur{}
\datecreate{2023-02-07}
\organisation{AMSCC}
\difficulte{}
\contenu{

\texte{ 	Un standard téléphonique est, en moyenne, occupé 6 fois sur 10. Soit $X$ le nombre de tentatives nécessaires pour obtenir le standard. }
	\begin{enumerate}
		\item \question{ Déterminer la loi de $X$. }
		\item \question{ Combien d'essais sont nécessaires en moyenne pour obtenir le standard ? }
		\item \question{ Quelle est la probabilité d'obtenir le standard en moins de 3 essais ? }
	\end{enumerate}}
